\documentclass[11pt]{article}
\usepackage[T1]{fontenc}
\usepackage{lmodern}
\usepackage[margin=1.08in]{geometry}
\usepackage{amsmath,amssymb,mathtools,microtype}
\usepackage{xurl}
\usepackage[colorlinks=true,linkcolor=blue,urlcolor=blue]{hyperref}
\hypersetup{pdftitle={The irrationality exponent of nonzero rational arctangents is 2},pdfauthor={Ryan Matthew Casper}}
\setlength{\emergencystretch}{3em}
\setcounter{secnumdepth}{0}
\providecommand{\tightlist}{\setlength{\itemsep}{0pt}\setlength{\parskip}{0pt}}
\title{The irrationality exponent of nonzero rational arctangents is 2}
\author{Ryan Matthew Casper}
\date{October 7, 2026}
\begin{document}
\maketitle
This manuscript presents an extension of OpenAI's weighted logarithmic
interpolation and determinant method to arctangents of nonzero rational
numbers. The arithmetic takes place in the Gaussian integers. The
argument retains the actual exponential factors in the matrix, the row
denominator savings, and the approximation-center translation costs.

The Lean endpoints prove the statement below without an irrationality,
non-torsion, interpolation, ampleness, or approximation-bound
hypothesis. The geometric comparison framework and its formal
infrastructure are inherited from OpenAI's released library and our
preceding rational-logarithm implementation. Substantial inherited
lemmas are identified explicitly below. This is a mathematical
exposition of the formal argument, not a claim of external human peer
review.

Yufei Bai's preprint submitted on 20 August 2026 reports the upper bound
8.585166\ldots{} for the irrationality measure of arctan(1/2). We cite
it as a concrete earlier benchmark; this investigation does not
establish an exhaustive priority claim.
\href{https://arxiv.org/abs/2608.19782}{Bai, arXiv:2608.19782}.

\section{1. Statement and conventions}\label{statement-and-conventions}

\textbf{Theorem.} For every rational number \(r\ne0\), the real number
\(\arctan r\) is irrational and \[
\mu(\arctan r)=2.
\] In particular, \(\mu(\arctan(1/2))=2\). The approximation statement
proved before taking the irrationality exponent is \[
\forall\nu>2\;\exists Q\in\mathbb Z,\ Q\ge2:\quad
\forall p,q\in\mathbb Z,\ q\ge Q\Longrightarrow
q^{-\nu}\le\left|\arctan r-\frac pq\right|.
\tag{1}
\] The threshold depends on \(r,\nu\), and is uniform in \(p,q\).
Fractions in (1) need not be reduced. The case \(r=0\) is excluded: its
arctangent is zero and exact rational approximations prevent (1).

For an irrational real number \(x\), \(\mu(x)\) is the supremum of the
positive exponents \(\nu\) for which infinitely many reduced rational
numbers \(p/q\), with \(q\ge2\), satisfy \(0<|x-p/q|<q^{-\nu}\). The
formal proof establishes irrationality and the lower bound 2 before
identifying the supremum, so an unbounded-supremum convention plays no
role.

We first prove a slightly broader angle statement. Let \(\xi\ne0\) be
real, let \(U\in\mathbb Z[i]\), and let \(D\) be a positive integer.
Suppose \[
\alpha=U/D=e^{i\xi},\qquad
\alpha^0,\alpha^1,\alpha^2,\ldots\text{ are pairwise distinct}.
\tag{2}
\] Then (1), with \(\xi\) in place of \(\arctan r\), holds. The
arctangent hypotheses and the passage from \(2\arctan r\) to
\(\arctan r\) are discharged in §§2 and 8.

\section{2. The arctangent exponential and the torsion
exceptions}\label{the-arctangent-exponential-and-the-torsion-exceptions}

For rational \(r\), put \[
\xi=2\arctan r,\quad \omega=i\xi,\quad
a_r=\frac{1-r^2}{1+r^2},\quad b_r=\frac{2r}{1+r^2}.
\] The double-angle formulas and \(\tan(\arctan r)=r\) give \[
e^{i\xi}=a_r+ib_r=\frac{1+ir}{1-ir},\qquad a_r^2+b_r^2=1.
\tag{3}
\] There is no choice of a global logarithm branch in this identity. We
specify the real angle \(2\arctan r\) directly. The local logarithm used
later is only \(\log(1+t)\) on \(|t|<1\).

Every number \(a_r+ib_r\) can be written \(U/D\) with Gaussian \(U\) and
positive integer \(D\). The family implementation uses \[
D=\operatorname{den}(a_r)\operatorname{den}(b_r),\qquad
U=\operatorname{num}(a_r)\operatorname{den}(b_r)
+i\operatorname{num}(b_r)\operatorname{den}(a_r).
\] Lowest terms are unnecessary. For the first target, the separate
formal path uses exactly \[
r=1/2,\quad \xi=2\arctan(1/2),\quad U=3+4i,\quad D=5.
\tag{4}
\]

We need the exact root-of-unity exceptions. Suppose \(z=a+ib\), with
\(a,b\in\mathbb Q\), \(a^2+b^2=1\), and \(z^n=1\) for a positive integer
\(n\). Both \(z\) and \(\bar z\) are algebraic integers. Hence \[
2a=z+\bar z,\qquad 2b=-i(z-\bar z)
\] are rational algebraic integers and therefore ordinary integers.
Write them as \(k,l\). Then \(k^2+l^2=4\), so \((k,l)\) is one of
\((\pm2,0),(0,\pm2)\). Thus \(z\in\{1,-1,i,-i\}\).

For (3), \(a_r=1\) forces \(r=0\); \(a_r=-1\) is impossible; and
\(a_r=0\) forces \(r=\pm1\). Conversely these exceptional values do give
roots of unity. Therefore the powers of \(\alpha\) are distinct whenever
\(r\notin\{0,1,-1\}\): an equality \(\alpha^j=\alpha^k\), with \(j<k\),
would give \(\alpha^{k-j}=1\). In particular the powers in (4) are
distinct. For that single value one can also use its trace
\(\alpha+\bar\alpha=6/5\), which cannot be an integer.

The nonzero rational exceptions \(r=\pm1\) will be handled separately
using \(\arctan(1)=\pi/4\), \(\arctan(-1)=-\pi/4\), and the certified
OpenAI bound for \(\pi\).

\section{3. Weighted logarithmic interpolation at distinct exponential
coordinates}\label{weighted-logarithmic-interpolation-at-distinct-exponential-coordinates}

We first describe the geometric input, because it is the point at which
the extension changes the original centers. The geometry and
interpolation are over \(\mathbb C\); the matrices used later have
Gaussian rational entries.

Let \(m,K\ge1\), and let \[
W=(w_0,\ldots,w_m),\qquad V=(v_0,\ldots,v_m)
\] be positive rational weights. Let \(\sigma>0\) be rational. Define
the conservative comparison constant \[
\mathcal C(m,\sigma)=
2(m+2)^{m+2}\left(1+\frac{m+2}{\sigma}\right)^{m+2}.
\] We impose the following conditions: \[
K(1+3\sigma)^{m+1}\frac{\prod_iw_i}{\prod_iv_i}<1,
\qquad
(1+\sigma)w_i<v_i\quad(1\le i\le m).
\tag{G1}
\] For equally sized subsets \(A,B\subseteq\{0,\ldots,m\}\), if the
largest index in their symmetric difference is positive and belongs to
\(A\), require \[
\mathcal C(m,\sigma)\prod_{j\in B}v_j
<\prod_{i\in A}w_i.
\tag{G2}
\]

The centers are \[
b_j=(y_j,c_{j1},\ldots,c_{jm}),\quad 0\le j<K,
\] with distinct nonzero \(y_j\). There is no requirement that the
\(c_{ji}\) be distinct. Near each center use formal coordinates \[
t=Y/y_j-1,\qquad
u_i=X_i-c_{ji}-\log(1+t).
\tag{G3}
\] A jet packet of height \(H\) consists of the coefficients of
\(t^su^\beta\) with \(v_0s+\sum_i v_i\beta_i<H\).

\textbf{Interpolation theorem used here.} There is a positive integer
\(R\), clearing the weight denominators, such that for every
sufficiently large integer \(n\), polynomials of \(W\)-weighted degree
at most \(H=nR\) map surjectively onto the packets of height \(H\) at
all \(K\) centers.

The cofinal heights \(nR\) suffice for the determinant argument. We do
not need, and do not claim as a checked intermediate theorem,
surjectivity for every sufficiently large real height.

\subsection{3.1 A curve estimate}\label{a-curve-estimate}

For an irreducible algebraic curve \(C\) in the affine space with
\(Y\ne0\), work on its smooth complete normalization. At every place
\(P\), set \[
D_W(C)=\sum_P
\max\left(0,\frac{\operatorname{pole}_P Y}{w_0},
\frac{\operatorname{pole}_P X_1}{w_1},\ldots,
\frac{\operatorname{pole}_P X_m}{w_m}\right).
\] The notation means the sum of the displayed local maxima; identically
zero coordinates contribute no poles. For each branch over a center
define \[
h_P=\min\left(\frac{\operatorname{ord}_P t}{v_0},
\frac{\operatorname{ord}_P u_1}{v_1},\ldots,
\frac{\operatorname{ord}_P u_m}{v_m}\right),\qquad
S(C)=\sum_{P\text{ over centers}}h_P.
\] Orders of identically zero series may be infinite. On a nonconstant
curve the minimum remains finite.

The needed estimate is \[
D_W(C)\ge(1+\sigma)S(C).
\tag{G4}
\]

Suppose instead that (G4) fails. The volume inequality in (G1) implies,
by counting monomials and jet conditions, that for every sufficiently
large \(N\) there is a nonzero polynomial \(F_N\) of \(W\)-degree at
most \(N\), whose \(V\)-weighted order at every center is at least
\((1+3\sigma)N\). This count needs no independence of the vanishing
conditions.

Use the commuting polynomial differential frame \[
D_0=Y\partial_Y+\sum_{i=1}^m\partial_{X_i},
\qquad D_i=\partial_{X_i}\quad(i>0).
\] In the chart (G3), these become \((1+t)\partial_t,\partial_{u_i}\).
Differentiation preserves the \(W\)-degree bound and loses at most its
assigned \(V\)-cost in weighted order.

Every derivative word of cost at most \(\sigma N\) therefore has
weighted order at least \((1+2\sigma)N\). Its restriction to each branch
has order at least \((1+2\sigma)Nh_P\), whereas its total pole degree is
at most \(ND_W(C)\). The assumed failure of (G4) forces all these
derivatives to vanish identically on \(C\).

Consider common zero loci of derivatives at successive cost cutoffs \[
r\delta_N,\qquad \delta_N=\sigma N/(m+2).
\] Choose nested irreducible components containing the affine part of
\(C\). Their dimensions range from 1 to \(m\); consecutive stages
eventually share a component \(Z\). On this component, the defining
equations from the earlier stage have degree at most \(N\), and their
further derivatives through cost \(\delta_N\) vanish.

We now use the inherited \textbf{persistent normal comparison lemma}.
For a coordinate normal basis indexed by \(A\), and a differential-frame
normal basis indexed by \(B\), persistence gives \[
\prod_{i\in A}w_i\le
\mathcal C(m,\sigma)\prod_{j\in B}v_j.
\tag{G5}
\] The lemma compares an upper bound from weighted Bézout with a lower
bound from rectangular transverse jets. Its large-\(N\) requirement is
uniform in the component: the relevant rectangular cutoffs are at least
2 once \(N\) is sufficiently large in terms of the fixed weights,
dimension and \(\sigma\).

This is a substantial reused lemma, not a new result about arbitrary
analytic frames. The formal applications are mapped in §9 to the
original persistent-weight and transverse-slice sources.

If some positive-index coordinate direction is not tangent to \(Z\),
take the largest such index \(i\). The coordinate normal basis can
include \(i\), with all larger directions tangent. Equations (G2) and
(G5) force every frame normal basis to include \(D_i\). All other frame
vectors lie in \[
\ker(dX_i-dY/Y).
\] Their normal rank is one less than the codimension, so the tangent
space of \(Z\) lies in this hyperplane. Thus \(dX_i=dY/Y\) on \(Z\).

If \(Y\) were nonconstant, restrict to an algebraic curve on which it
remains nonconstant and pass to a smooth complete model. An exact
differential \(dX_i\) has zero residues, while \(dY/Y\) has residue
\(\operatorname{ord}_P Y\) at each zero or pole. A nonconstant rational
function has such a place, giving a contradiction. Consequently \(Y\) is
constant on \(Z\).

In the remaining case all positive-index coordinate directions are
tangent. Since \(Z\) is proper, its tangent space is their span;
therefore \(dY=0\), and characteristic zero again makes \(Y\) constant.
In either case \(Y\) is constant on \(C\).

Distinct \(y_j\) now matter: a constant-\(Y\) curve meets at most one
center. Some \(X_i\) is nonconstant on \(C\). At that center \(t=0\) and
\(u_i=X_i-c_{ji}\), so the zero/pole equality for this function gives \[
v_iS(C)\le w_iD_W(C).
\] The coordinate inequality in (G1) contradicts the assumed failure of
(G4). This proves the curve estimate, counting every branch, including
branches over singular points.

\subsection{3.2 From curves to bounded-degree
interpolation}\label{from-curves-to-bounded-degree-interpolation}

Choose \(R\) so that \(R/w_i,R/v_i\) are positive integers. Compactify
by the monomials of \(W\)-degree at most \(R\), retaining the constant
coordinate as the affine chart. If \(L=\mathcal O(1)\), then on every
normalized affine curve, \[
\deg L=R D_W(C).
\] The constant and pure-coordinate monomials attain the local maxima
defining this degree.

Replace the logarithm in (G3) by a polynomial truncation whose first
omitted term has \(t\)-weight larger than every relevant \(v_i\). This
leaves the weighted filtration and branch contacts unchanged. At center
\(j\), form the finite-colength ideal \[
I_j=(t^{R/v_0},(u'_1)^{R/v_1},\ldots,(u'_m)^{R/v_m}).
\] These ideals give a coherent ideal \(I\), equal to the unit ideal
away from the centers.

On the ordinary blowup, write \(I\mathcal O=\mathcal O(-E)\) and
\(A=\pi^*L\). The source constructs the actual curve models and proves,
for noncontracted affine curves, \[
A\cdot C=RD_W(C),\qquad E\cdot C=RS(C).
\] Affine curves with \(Y\equiv0\) meet none of the prescribed centers,
so \(S(C)=0\) and their curve estimate is automatic. Together with (G4)
for curves meeting \(Y\ne0\), this makes \(A-(1+\sigma)E\) nef. Boundary
curves avoid \(E\); on contracted curves \(A\) has degree zero and
\(-E\) is positive by relative ampleness.

For a sufficiently large integer \(r>1\), \(rA-E\) is ample. The exact
identity \[
A-E=
\frac{r-1}{r(1+\sigma)-1}\,[A-(1+\sigma)E]
+\frac{\sigma}{r(1+\sigma)-1}\,[rA-E]
\tag{G6}
\] has positive coefficients, and supplies ampleness together with a
uniform positive curve margin. This step does not infer ampleness from
strict nefness alone.

For all sufficiently large \(n\), the ordinary blowup has \[
\pi_*\mathcal O(-nE)=I^n,\qquad R^q\pi_*\mathcal O(-nE)=0\quad(q>0).
\] These are eventual identities for ordinary powers of \(I\); they are
not assertions about integral closures or about every small \(n\). Serre
vanishing for \(A-E\), followed by the projection formula and Leray,
yields \[
H^1(X,I^nL^n)=0.
\] Sections therefore surject onto the local quotients. The monomials
generating \(I_j^n\) have weighted order at least \(nR\), so this
surjection prescribes all packets below \(H=nR\). Completion does not
change these finite-colength quotients.

Finally, eventual polynomial generation for the projective embedding
expresses sections of \(L^n\) as degree-\(n\) homogeneous polynomials in
its coordinates. On the affine chart their \(W\)-degree is at most
\(nR\). This last step is essential: a Chinese remainder theorem without
a degree bound would not give our interpolation statement.

\section{4. The actual matrix}\label{the-actual-matrix}

Fix the data in (2) and an exponent \(\nu>2\). Assume for contradiction
that there are approximations \(p/q\) with arbitrarily large positive
integer denominators and \[
|\xi-p/q|\le q^{-\nu}.
\] Choose finitely many of them, denoted \(r_i=p_i/q_i\), whose
denominators will be taken very large and successively separated in §7.
Set \[
w_i=\lceil\log q_i\rceil,\qquad v_i=w_i/\theta\quad(1\le i\le m),
\] with rational \(0<\theta<1\). Use the multiplicative centers
\(y_j=\alpha^j\) and additive centers \(c_{ji}=ijr_i\), for
\(0\le j<K\). Here the isolated symbol \(i\) before \(jr_i\) denotes the
imaginary unit; subscripts index the chosen approximants. Equivalently,
the additive center in coordinate \(k\) is \(\sqrt{-1}\,j p_k/q_k\).

Choose \(F_0>2/\theta\), and define \[
T_i=\lceil F_0w_i/v_0\rceil,\qquad
G_i(t)=\sum_{1\le k<T_i}\frac{(-1)^{k+1}t^k}{k}.
\] The omitted tail starts at order \(T_i\), whose weight satisfies
\(T_iv_0\ge F_0w_i>v_i\). It therefore preserves the jet packets in §3.

Columns are monomials \(Y^hX^\gamma\) with \(w_0h+w\gamma\le H\). Rows
are triples \((j,s,\beta)\) with \[
0\le j<K,\qquad v_0s+w\beta/\theta<H.
\] The entry in the actual matrix is \[
[t^su^\beta]\,\bigl(\alpha^j(1+t)\bigr)^h
\prod_i\bigl(\sqrt{-1}\,jr_i+G_i(t)+u_i\bigr)^{\gamma_i}.
\tag{5}
\] The exponential factor \(\alpha^{jh}\) is retained. All entries lie
in \(\mathbb Q(i)\).

The interpolation theorem supplies full row rank at cofinal heights
\(H=nR\). Select a square nonzero minor \(\Delta_H\) containing every
row. Let \(M_H\) be the number of rows, and write \[
b_H=\frac{\sum_{\text{rows}}w\beta}{M_HH},\qquad 0\le b_H\le\theta.
\] Weighted lattice counting gives \[
M_H\sim\frac{K\theta^mH^{m+1}}{(m+1)!v_0\prod_iw_i}.
\tag{6}
\] Both determinant bounds below concern this same minor, and are
uniform in the selected columns.

\section{5. Gaussian arithmetic and the row
savings}\label{gaussian-arithmetic-and-the-row-savings}

Let \[
L_i=\operatorname{lcm}(1,\ldots,T_i),\qquad
E_H=\prod_iL_i^{\lfloor H/w_i\rfloor}.
\] Multiply column \((h,\gamma)\) by \(D^{Kh}\prod_iq_i^{\gamma_i}\),
row \((j,s,\beta)\) by \(\prod_iq_i^{-\beta_i}\), and every entry by
\(E_H\). The row factors are essential; discarding them loses the
leading saving in (8).

When \(\beta\le\gamma\), the cleared entry is \[
E_H U^{jh}D^{(K-j)h}\binom{\gamma}{\beta}
[t^s](1+t)^h
\prod_i\bigl(\sqrt{-1}\,jp_i+q_iG_i(t)\bigr)^{\gamma_i-\beta_i}.
\tag{7}
\] Otherwise it is zero. The identity
\(D^{Kh}\alpha^{jh}=U^{jh}D^{(K-j)h}\) uses \(j\le K\). Each
\(\gamma_i-\beta_i\le\lfloor H/w_i\rfloor\), so the displayed lcm powers
clear all coefficient denominators. Thus the cleared entries, and their
determinant, lie in \(\mathbb Z[i]\).

All scaling factors are nonzero. Since \(\Delta_H\ne0\), its cleared
determinant is a nonzero Gaussian integer \(Z\). Consequently \(|Z|^2\)
is a positive integer and \(|Z|\ge1\). This does not introduce a
degree-two multiplier: the two complex embeddings give the same absolute
value, and the lower bound is already in the chosen complex absolute
value.

Put \(\Lambda=4+\log4\) and \(w_* =\min_iw_i\). The inherited
conservative lcm estimate is \(\log L_i\le\Lambda T_i\). Column degree,
\(T_i\le F_0w_i/v_0+1\), and \(w_i-1\le\log q_i\le w_i\) give \[
\frac{\log|\Delta_H|}{M_HH}\ge -(1-b_H)-E_{\rm ar},
\tag{8}
\] where \[
E_{\rm ar}=\frac{\Lambda F_0m}{v_0}
+\Lambda\sum_i\frac1{w_i}+\frac{\theta}{w_*}
+\frac{K\log D}{w_0}.
\tag{9}
\] Indeed, the common clearing factor costs at most the first two terms.
The selected columns cost at most \(M_HH(1+K\log D/w_0)\). The row
logarithms save at least \(M_HH(b_H-\theta/w_*)\). Combining these three
inequalities yields (8). The Gaussian numerator needs no further
clearing factor and contributes no additional lower-bound cost.

\section{6. Complex analytic expansion and collision
saving}\label{complex-analytic-expansion-and-collision-saving}

Set \(\omega=i\xi\) and \(\rho=100\max(1,|\xi|)\). For a column monomial
\(P=Y^hX^\gamma\) and multi-index \(b\), define the entire function \[
f_{b,P}(z)=[u^b]P(e^z,z+u_1,\ldots,z+u_m)
=\binom{\gamma}{b}e^{hz}z^{|\gamma|-|b|},
\] with value zero if \(b\not\le\gamma\). On \(|t|<1\), put \[
z_j(t)=j\omega+\log(1+t).
\] Since \(e^{j\omega}=\alpha^j\), we have \(e^{z_j(t)}=\alpha^j(1+t)\).
The additive arguments in (5) become \[
z_j(t)+u_i+e_{ji}+\tau_i(t),\quad
e_{ji}=\sqrt{-1}\,j(r_i-\xi),\quad
\tau_i=G_i-\log(1+t).
\tag{10}
\] No step replaces \(j\omega\) by the principal value of
\(\log(\alpha^j)\). Such a replacement would be invalid when the powers
wind around the unit circle; it is absent from both the argument and the
formalization.

Expanding (10) gives a finite linear combination of row vectors \[
\bigl([t^\ell]f_{b,P}(z_j(t))\bigr)_P,
\] with \(b\ge\beta\), \(wb\le H\), and coefficients independent of the
selected column. The rotation has modulus one, so the approximation
error satisfies \[
|e_{ji}|=j|r_i-\xi|\le 2K e^\nu e^{-\nu w_i}.
\] Each tail starts at \(T_i\). A term using tail multiplicities \(e_i\)
can contribute only if \(\sum_i e_iT_i\le s\), whence
\(\sum_iw_ie_i<H/F_0\). On the circle \(|t|=1/2\), the inherited
coefficient and binomial estimates bound its scalar by \[
\exp\{-\nu w(b-\beta)+HE_{\rm tr}\},\qquad
E_{\rm tr}=\frac\nu{F_0}+\frac{\log2}{v_0}
+\frac{\log4+\log(2K)+\nu}{w_*}.
\tag{11}
\] This accounts for the approximate additive centers; replacing them by
exact exponential centers would omit this term.

On \(|z|\le\rho K\), the entire functions are bounded by \[
\exp\left\{H\left(\frac{\rho K}{w_0}
+\frac{\log(2\rho K)}{w_*}\right)\right\}.
\] The center estimate \(|j\omega|+3/4\le\rho K/2\) places \(z_j(t)\)
within the smaller disc when \(|t|=1/2\). Taylor expansion in
\(z/(\rho K)\) is therefore absolutely convergent with a geometric
majorant. If two rows use the same transverse index \(b\) and the same
Taylor order, that determinant term is zero. Among \(n_b\) rows with a
common \(b\), the Taylor orders of a surviving term have sum at least
\(\binom{n_b}{2}\).

The inherited collision determinant estimate, with conservative constant
\(c=\log2/4\), gives \[
|\det B|\le\exp\left\{-c\sum_bn_b^2
+M_HH\bigl(E_{\rm hol}+o(1)\bigr)\right\},
\] where \[
E_{\rm hol}=\frac{\rho K}{w_0}+\frac{\log2}{v_0}
+\frac{\log(2\rho K)}{w_*},\qquad E_{\rm an}=E_{\rm tr}+E_{\rm hol}.
\tag{12}
\] The estimate also holds for singular test matrices. Its proof expands
the determinant of the entire-function row tests, uses the
repeated-order cancellation just described, and bounds the remaining
absolutely convergent geometric sums and factorial factors. The formal
source applies the existing theorem for arbitrary complex \(\omega\),
retaining the factor \(\exp(j\omega)^h=\alpha^{jh}\).

Choose \(A>\theta\) and \(0<\eta<1\). Let \(N_A(H)\) count multi-indices
with \(wb\le AH\), and set \[
c_H=\frac{c\eta^2M_H}{H N_A(H)}.
\] If at least \(\eta M_H\) rows use those low indices, Cauchy--Schwarz
gives \(\sum_bn_b^2\ge\eta^2M_H^2/N_A(H)\), hence the normalized saving
\(-c_H\). Otherwise more than \((1-\eta)M_H\) rows have \(wb>AH\), and
the product of scalar factors in (11) gives the saving
\(-\nu(A(1-\eta)-b_H)\).

For fixed parameters, the finite expansion has at most \[
Q_H=(\lfloor H\rfloor+1)^{2m}(\lfloor H/v_0\rfloor+1)
\] choices per row. Determinant multilinearity introduces at most
\(Q_H^{M_H}\) choices, adding only \(\log Q_H/H=o(1)\) after
normalization. Together with the normalized factorial remainder, this
proves \[
\frac{\log|\Delta_H|}{M_HH}
\le E_{\rm an}+o(1)+
\max\{-c_H,-\nu(A(1-\eta)-b_H)\}.
\tag{13}
\] All parameters and approximations are fixed before \(H\) tends to
infinity. Lattice counting also gives \[
c_H\longrightarrow c_\infty=
\frac{c\eta^2K\theta^m}{(m+1)v_0A^m}.
\tag{14}
\]

\section{7. Parameter choice and
contradiction}\label{parameter-choice-and-contradiction}

Fix \(\xi,U,D,\nu\) first. The quantities \(\rho\) and \(\delta=\log D\)
are fixed. Choose positive rational \(\theta,A,B,C\) such that \[
0<\theta<A<B<1,\quad
\nu(A-\theta)>1-\theta,\quad
C>1,\quad CB<1,\quad B<C\theta<1.
\tag{15}
\] These inequalities are compatible precisely with the required range
\(\nu>2\) in this construction: choose rational \(b\in(1/2,1-1/\nu)\)
and small rational \(s>0\), then set \(\theta=1-s\), \(A=1-bs\). The
approximation gap is positive and \(A^2<\theta\) for sufficiently small
\(s\). Take \(C\in(A/\theta,1/A)\), then \(B\in(A,\min(1/C,C\theta))\).

Take small rational \(\eta>0\) with \[
g=\nu(A(1-\eta)-\theta)-(1-\theta)>0,
\qquad \epsilon=\min(g/2,1/2).
\] Choose \(F_0>2/\theta\) with \(\nu/F_0<\epsilon/3\). Now increase
\(m\), setting \[
K=\lfloor C^m\rfloor,\quad w_0=B^{-m},\quad
v_0=2K\theta^mw_0.
\] The ambient interpolation mass is \(1/2\), while \[
K/w_0\le(CB)^m\to0,\quad
v_0\sim2(C\theta/B)^m\to\infty,\quad m/v_0\to0,
\] and \[
c_\infty=\frac{c\eta^2(B/A)^m}{2(m+1)}\to\infty.
\] Choose one finite \(m\ge1\) for which \[
\frac{\Lambda F_0m+2\log2}{v_0}
+\frac{(\rho+\delta)K}{w_0}<\epsilon/3,
\qquad c_\infty>2.
\tag{16}
\] Then fix these parameters. Since \(K\theta^m<1\), choose a small
positive rational \(\sigma\) with \[
(1+3\sigma)^{m+1}/2<1,\quad
(1+3\sigma)^mK\theta^m<1,\quad
(1+\sigma)\theta<1.
\] Only now choose the finitely many approximations. Arbitrarily large
denominators permit all weights to be large and successively separated,
meeting the product comparisons in §3 and \[
\Lambda\sum_i\frac1{w_i}+
\frac{\theta+\log4+\log(2K)+\nu+\log(2\rho K)}{w_*}<\epsilon/3.
\tag{17}
\] There are finitely many product comparisons. Common factors cancel
because \(v_i=w_i/\theta\) for positive indices, and increasing the next
denominator meets each remaining comparison. The interpolation threshold
may depend on the chosen centers; that causes no circularity, since the
centers are fixed before the final limit in \(H\).

Equations (9), (11), (12), (16), and (17) imply \[
E_{\rm ar}+E_{\rm an}<\epsilon<g,\qquad
c_\infty>1+E_{\rm ar}+E_{\rm an}.
\tag{18}
\] In particular the denominator cost \(K\log D/w_0\) and complex radius
cost \(\rho K/w_0\) both tend to zero at the dimension-selection stage.
There is no unaccounted conjugate term.

For sufficiently large cofinal \(H=nR\), choose the nonzero minor
furnished by interpolation. If the maximum in (13) is the collision
alternative, compatibility with (8) requires \[
c_H\le1-b_H+E_{\rm ar}+E_{\rm an}+o(1)
\le1+E_{\rm ar}+E_{\rm an}+o(1),
\] contradicting (14) and (18). In the other alternative, compatibility
requires \[
\nu(A(1-\eta)-b_H)-(1-b_H)
\le E_{\rm ar}+E_{\rm an}+o(1).
\] The left side is at least \(g\), since it decreases with \(b_H\) and
\(b_H\le\theta\). This contradicts (18) as well. Thus arbitrarily large
approximations of the assumed quality cannot exist, proving the angle
statement (2).

\section{8. Scaling, the exceptional cases, and the
exponent}\label{scaling-the-exceptional-cases-and-the-exponent}

We use an explicit transport of the approximation bound. Suppose \(x\)
satisfies (1) and \(n\) is a positive integer. Fix \(\nu>2\), and choose
\(\mu=(\nu+2)/2\), so \(2<\mu<\nu\). For sufficiently large \(q\), the
bound for \(x\) and \(q^{\nu-\mu}\ge n\) give \[
\left|\frac xn-\frac pq\right|
=\frac1n\left|x-\frac{np}{q}\right|
\ge\frac1n q^{-\mu}\ge q^{-\nu}.
\tag{19}
\] Negation preserves the same statement by replacing \(p\) with \(-p\).
These are formal lemmas on the fully quantified eventual bound; no
invariance under algebraic scaling is presumed.

For \(r\notin\{0,1,-1\}\), §2 gives the data (2) for \(x=2\arctan r\),
and (19) with \(n=2\) proves (1). The dedicated \(r=1/2\) path uses (4).
For \(r=\pm1\), apply (19) with \(x=\pi,n=4\), using OpenAI's certified
\(\pi\) bound, and use negation for \(-1\). This proves (1) for every
rational \(r\ne0\).

Finally, if \(\arctan r=u/v\) were rational, the unreduced fractions
\(ku/(kv)\) would violate (1) at arbitrarily large denominators. Hence
it is irrational. Dirichlet approximation supplies infinitely many
reduced approximants of exponent 2, while (1) excludes every larger
exponent. Therefore its irrationality exponent is exactly 2.

\section{9. Formal boundary, verification, and
attribution}\label{formal-boundary-verification-and-attribution}

The endpoint source is
\href{lean/Arctangent/Main.lean}{Arctangent/Main.lean}. The internal
generic angle theorem retains the data in (2), but both public
arctangent endpoints discharge them. The family endpoint has only the
hypothesis \(r\ne0\).

\begin{itemize}
\tightlist
\item
  Specified angle and Gaussian rational base:
  \href{lean/Arctangent/PeriodData.lean}{PeriodData}
\item
  Rational-coordinate roots of unity:
  \href{lean/Arctangent/GaussianRoots.lean}{GaussianRoots}
\item
  Double-angle exponential, all rational inputs:
  \href{lean/Arctangent/ArctanCenters.lean}{ArctanCenters}
\item
  Distinct-center interpolation geometry:
  \href{lean/Logarithm/Interpolation.lean}{Inherited Interpolation},
  \href{lean/Arctangent/GlobalMatrixInterpolation.lean}{GlobalMatrixInterpolation}
\item
  Gaussian clearing and exact row savings:
  \href{lean/Arctangent/ScaledArithmetic.lean}{ScaledArithmetic}
\item
  Rotated approximation errors:
  \href{lean/Arctangent/ScaledApproximationErrors.lean}{ScaledApproximationErrors}
\item
  Complex collision estimate:
  \href{lean/Logarithm/ScaledCollision.lean}{Inherited ScaledCollision},
  \href{lean/Arctangent/LiteralAnalyticSummand.lean}{LiteralAnalyticSummand}
\item
  Full determinant expansion and upper bound:
  \href{lean/Arctangent/AnalyticAggregate.lean}{AnalyticAggregate}
\item
  All parameter costs and final contradiction:
  \href{lean/Logarithm/AdmissibleParameters.lean}{Inherited
  AdmissibleParameters},
  \href{lean/Arctangent/Parameters.lean}{Parameters},
  \href{lean/Arctangent/DeterminantContradiction.lean}{DeterminantContradiction}
\item
  Division and negation transport:
  \href{lean/Arctangent/Scaling.lean}{Scaling}
\item
  Inherited normal comparison:
  \href{lean/OAI/NumberTheory/PiExponent/Approximation/PersistentWeightComparison.lean}{PersistentWeightComparison},
  \href{lean/OAI/NumberTheory/PiExponent/Jets/TransverseSliceComparison.lean}{TransverseSliceComparison}
\end{itemize}

The new Lean sources are separate from the preceding audited
rational-logarithm bundle. The source manifest, compiler and checker
pins, reproduction commands, hashes, verification results, and negative
controls are summarized in \href{lean/VERIFICATION.md}{VERIFICATION.md}.
Those records distinguish the initial Lean build, transitive axiom
audits, independent exported-proof checking, and replay through a second
stock Lean kernel. They do not claim that a human referee has
reconstructed the proof or that the priority search is exhaustive.

The mathematical framework and original library are credited to
\href{https://github.com/openai/math/tree/adc7f1241b42e322a6451854ab7e4b4c146bf78a}{OpenAI/math,
commit adc7f1241b42e322a6451854ab7e4b4c146bf78a}. Our preceding
positive-rational-logarithm package supplies the distinct-center
geometric generalization and generic complex analytic interfaces used
here. We make no exhaustive priority or novelty claim.

\textbf{AI use disclosure.} Ryan Matthew Casper initiated and directed
the investigation. OpenAI coding agents contributed substantially to
assessing the Gaussian extension, developing the mathematical argument,
adapting and writing the Lean proofs, running verification and automated
checks, and drafting this exposition. The work reuses OpenAI's method
and released library, Mathlib, and the author's preceding AI-assisted
rational-logarithm formalization. Automated review and kernel-checking
results are evidence about the specified formal statements; they are not
external human peer review. No external referee was contacted in
preparing this package.

The central new checks are the exact Gaussian clearing identity (7),
preservation of the row savings, the imaginary approximation-center
error (10), the retained exponential factor in the analytic matrix, all
fixed-base costs in (16), and the explicit treatment of torsion and
scaling. The inherited geometric comparison and ordinary-power
interpolation remain substantial parts of the argument and should be
included in any mathematical review.
\end{document}
